% =====================================================================
% Explainable AI-Based Personalized Learning Path Recommendation System
% for University Students  --  REVISED
%
% Your paper, with fixes applied. Section order, subsection titles,
% references and most sentences are unchanged. What changed:
%   - every fabricated number replaced with a measured one
%   - DQN removed (never implemented); planner described as what it is
%   - user study removed; objective explanation metrics in its place
%   - one new results subsection (V-E): the cold-concept finding
%   - Limitations + Ethics added to VI
%   - factual fixes (drug diversion, "Dynamic" DKT, 30-min sessions,
%     hardware, graph size, libraries that were not used)
%
% FIGURES: upload fig1.png fig2.png fig3.png fig4.png. If a file is
% missing the paper still compiles and shows a placeholder box, so the
% build never breaks. Change the filenames in the \figslot lines if
% yours differ.
%
% Two equations only (mining rule, forgetting decay). To remove them,
% delete each \begin{equation} block and the sentence introducing it.
%
% 8 pages with all four figures. Compiles clean under pdfLaTeX.
% =====================================================================

\documentclass[conference]{IEEEtran}
\IEEEoverridecommandlockouts

\usepackage{cite}
\usepackage{amsmath,amssymb,amsfonts}
\usepackage{graphicx}
\usepackage{textcomp}

% Figure slot: draws your image if the file is present, a labelled placeholder
% if it is not, so the paper compiles either way. Upload fig1.png ... fig4.png
% (or change the filenames below to match your own).
\newcommand{\figslot}[3]{%
  \begin{figure}[t]\centering
  \IfFileExists{#1}{\includegraphics[width=\columnwidth]{#1}}%
    {\framebox[\columnwidth]{\parbox[c][3.2cm][c]{0.9\columnwidth}{\centering\footnotesize\textit{#1 not found. Upload it.}}}}%
  \caption{#2}\label{#3}\end{figure}}

\begin{document}

\title{Explainable AI-Based Personalized Learning Path\\
Recommendation System for University Students}

\author{
\IEEEauthorblockN{Abhay Pratap Singh}
\IEEEauthorblockA{\textit{Department of Computer Science}\\
\textit{Chandigarh University}\\
Mohali, India\\
23BCS11784@cuchd.in}
\and
\IEEEauthorblockN{Akshay Guleria}
\IEEEauthorblockA{\textit{Department of Computer Science}\\
\textit{Chandigarh University}\\
Mohali, India\\
23BCS10517@cuchd.in}
\and
\IEEEauthorblockN{Uday Chaturvedi}
\IEEEauthorblockA{\textit{Department of Computer Science}\\
\textit{Chandigarh University}\\
Mohali, India\\
23BCS10807@cuchd.in}
\and
\IEEEauthorblockN{Gursimar Makkar}
\IEEEauthorblockA{\textit{Department of Computer Science}\\
\textit{Chandigarh University}\\
Mohali, India\\
23BCS11455@cuchd.in}
\and
\IEEEauthorblockN{Shikha Kamal}
\IEEEauthorblockA{\textit{Department of Computer Science}\\
\textit{Chandigarh University}\\
Mohali, India\\
shikha.e12552@cumail.in}
}

\maketitle

% =====================================================================
\begin{abstract}
As digital learning platforms have become increasingly common in higher education, they
have generated huge amounts of data on student interactions and performance that enable
personalised learning route recommendations. But present educational recommender systems
usually provide little transparency, cannot accurately simulate the time-varying
forgetfulness and knowledge of students, and are rarely evaluated in a way that
establishes which of their components actually contribute. This paper proposes an
Explainable Artificial Intelligence (XAI)-based personalized learning path recommendation
system consisting of a knowledge-graph-enhanced Deep Knowledge Tracing module, a
prerequisite-constrained planner, and a priority-backtracking explanation module.
Experiments on the Open University Learning Analytics Dataset (OULAD), under five-fold
cross-validation with folds grouped by student, give an AUC-ROC of $0.9538 \pm 0.0019$ and
an RMSE of $0.2621 \pm 0.0026$, exceeding DKT ($0.9435$, $p = 0.0009$) and a GNN baseline
($0.9502$, $p = 0.022$) after Holm correction. A central finding concerns evaluation
rather than architecture. Removing knowledge-graph propagation changes accuracy by
$-0.0001$ ($p = 0.77$), which would ordinarily be read as the graph contributing nothing;
that reading is an artefact of the measurement, because accuracy is scored only at
assessment positions and only 89 of 237 concepts carry assessments. Under a cold-concept
protocol that withholds all assessments for a subset of concepts and scores only those,
the same graph improves accuracy by $+0.0325$ ($p = 0.0003$, Cohen's $d = 5.06$), winning
on every fold. The planner produces zero prerequisite violations across 906 decisions
evaluated against the concepts learners actually studied next, and explanation quality is
assessed by objective fidelity, sufficiency and stability rather than by subjective
rating. All code and results are released publicly.
\end{abstract}

\begin{IEEEkeywords}
Explainable AI (XAI), personalized learning path, knowledge tracing, knowledge graphs,
educational data mining, ablation methodology, OULAD dataset.
\end{IEEEkeywords}

% =====================================================================
\section{Introduction}
\label{sec:intro}

Higher education's quick embrace of digital learning systems has opened up new
possibilities for individualized, data-driven instruction. Yet despite differences in
prior knowledge, learning speed and career aspirations, the learning process remains the
same for all students in conventional courses. In complex domains such ``one size fits
all'' approaches cause course mis-selection, knowledge gaps and less effective learning
\cite{li2026survey}--\cite{kuzilek2017oulad}. Personalized learning path recommendation
systems can overcome these constraints by evaluating student interactions to determine
appropriate activities, but most systems in use concentrate on prediction accuracy and
offer little explanation for any given recommendation \cite{corbett1994knowledge}.

Designing a system that describes a learner's evolving knowledge state, sequences
activities optimally, and explains its choices is therefore a major challenge. Most
knowledge tracing methods assess learning rather than optimize progression, while
collaborative filtering suffers from sparsity and weak personalization
\cite{piech2015dkt, zhang2017dkvmn}. Automated sequencing can optimize for long-term
outcomes \cite{ghosh2020akt}\cite{shen2021lpkt}, but its decisions remain opaque to
students and advisers, and the connection between knowledge graphs, sequential
decision-making and temporal knowledge modeling is still underexplored in education
\cite{kipf2017gcn, mnih2015dqn}.

A further limitation motivates much of this work. Such systems combine several
components --- a sequence model, a structural prior, a planner --- each justified by an
ablation showing that its removal degrades accuracy. We find this evaluation can be
systematically uninformative: when a structural prior exists precisely to support concepts
whose supervision is scarce, and accuracy is measured only where supervision is abundant,
the ablation tests the component in the one regime where it cannot matter.

The main contributions of this work are as follows.

\begin{enumerate}
\item A knowledge-graph-enhanced knowledge tracing model with a learned per-head temporal
decay, attaining AUC-ROC $0.9538 \pm 0.0019$ and calibration error $0.0054$ on OULAD under
student-grouped five-fold cross-validation.

\item Evidence that conventional ablation cannot detect the contribution of a knowledge
graph, together with a cold-concept protocol that can. Prerequisite propagation changes
standard accuracy by $-0.0001$ ($p = 0.77$) yet improves accuracy on concepts lacking
direct supervision by $+0.0325$ ($p = 0.0003$, $d = 5.06$).

\item A prerequisite-constrained planner whose pedagogical guarantee is enforced
structurally rather than learned, verified at zero violations across 906 recommendation
decisions evaluated against observed learner behaviour.

\item Three objective measures of explanation quality --- fidelity, sufficiency and
stability --- in place of subjective rating. One revealed that explanations accounted for
only $0.2\%$ of the benefit they described, a defect no rating scale would surface.

\item A reproducible pipeline from raw data to every reported figure, released publicly.
\end{enumerate}

% =====================================================================
\section{Literature Review}
\label{sec:lit}

Educational recommender systems have progressed from one-size-fits-all delivery to
adaptive personalization. Early content-based and collaborative filtering approaches
suffered cold-start problems, data sparsity and limited scalability
\cite{li2026survey}\cite{aher2013combination}. Hybrid models improved accuracy and
robustness \cite{kuzilek2017oulad} but forecast learner preferences rather than modeling
knowledge acquisition or optimizing sequential pathways \cite{corbett1994knowledge}.

Deep Knowledge Tracing (DKT) \cite{piech2015dkt} simulates the latent knowledge state of
learners over time. Unlike Bayesian Knowledge Tracing, which uses static parameters, DKT
learns temporal dependencies directly from data \cite{zhang2017dkvmn}, and attention-based
successors \cite{pandey2019sakt, ghosh2020akt} attend directly to any earlier interaction
rather than propagating through a recurrent chain \cite{vaswani2017attention}.
Reinforcement Learning enhances path recommendation by formalizing the process as a Markov
Decision Process and maximizing long-term learning rather than engagement; combined with
knowledge tracing it has been reported to reduce dropout and completion time
\cite{shen2021lpkt, zhao2026rlkt}. Knowledge graphs offer a complementary route to
structured, explainable recommendation by explicitly modelling concepts, course units and
their interactions \cite{kipf2017gcn, zhang2023finegrained}, with Bloom's taxonomy
grounding path design \cite{khosravi2022xaied} and graph neural networks handling the
resulting relational structure. Multi-objective frameworks have optimized dropout risk,
satisfaction and completion \cite{lundberg2017shap}\cite{ribeiro2016lime}, and
neuro-symbolic and generative approaches have begun to explore explanation generation
\cite{zhang2023finegrained, yang2025explainable}.

Four gaps remain. First, explainability is added as an afterthought rather than being
integral to recommendation optimization \cite{xie2025profile}; SHAP-based explanations
dominate in education and are applied mainly at the global level \cite{zhao2026rlkt}, and
where explanations are evaluated at all it is usually by participant rating, which measures
whether readers found an explanation agreeable rather than whether it is faithful to the
model. Second, many frameworks ignore the dynamic evolution of knowledge states, including
forgetting, compromising the pedagogical integrity of suggested trajectories
\cite{boujmiraz2026review}. Third, applying knowledge graphs to sequential decision-making
and explanation remains at an early stage \cite{qiang2026xaiedu}. Fourth, and most relevant
here, the contribution of such structural priors is customarily established by an ablation
on aggregate accuracy, without examining whether that metric can observe the component at
all.

\subsection{Research Gap}

Existing educational recommender systems focus on prediction and preference-based
suggestion, rarely integrating knowledge graphs, temporal knowledge tracing, constrained
path optimization and pedagogically meaningful explanation. This limits customization,
transparency and effective learning-path optimization. A further gap concerns validation:
component contributions are reported on metrics that may be structurally unable to detect
them, and explanation quality is reported by opinion rather than measurement.

% =====================================================================
\section{Methodology}
\label{sec:method}

\figslot{fig1.png}{Overall XAI-based personalized learning recommendation framework}{fig:framework}

\subsection{Overall Framework}

The proposed system combines Knowledge Graph Modeling, Explainable AI (XAI),
prerequisite-constrained path optimization and Deep Knowledge Tracing to recommend
personalized learning trajectories, through learner knowledge estimation, sequential path
optimization and explanation production \cite{piech2015dkt}. Student interaction history is
first processed to estimate the learner's changing knowledge state; prerequisite and
co-requisite relationships in the knowledge graph then constrain a planner that selects
the next learning task; finally each recommendation receives a pedagogical explanation
from a priority-backtracking system \cite{zhang2023finegrained}.

\subsection{Data Acquisition and Pre-processing}

Evaluation uses the Open University Learning Analytics Dataset (OULAD)
\cite{kuzilek2017oulad}, covering clicks, enrollments, assessment results and student
attributes. Interaction records are ordered chronologically, categorical features encoded
and numerical attributes standardized. OULAD timestamps are day-granular with no intra-day
resolution, so a session is one learner-day rather than an inactivity threshold. Temporal
features capture session volume, elapsed time since the previous interaction, and mastery
change. The processed data is split 80:10:10.

OULAD supplies no concept annotations, prerequisite relations or material titles, so the
concept graph is derived rather than read off. A concept is a (module, week-of-study)
pair: the material a module presents in a given week. Week placement uses the
\texttt{week\_from} field where present ($17.6\%$ of materials) and otherwise the median
access day; on the 1{,}121 materials carrying both, the two agree at Spearman
$\rho = 0.975$. Edges come from three rules: sequence edges between consecutive weeks,
assessment-coverage edges from the four preceding weeks to an assessment's week, and edges
mined from behaviour. A mined edge between concepts $a$ and $b$ is accepted only where
behaviour shows both consistent ordering and a measurable effect on later success,
%
\begin{equation}
\begin{split}
\mathrm{prec}(a,b) &= \frac{\left|\{u : \tau_u(a) < \tau_u(b)\}\right|}{\left|\{u : a, b \in h_u\}\right|},\\[2pt]
\mathrm{lift}(a,b) &= \frac{P(y_b \mid a \prec b)}{P(y_b)},
\end{split}
\label{eq:mining}
\end{equation}
%
where $\tau_u(\cdot)$ is learner $u$'s first interaction with a concept, requiring support
above 100 learners, $\mathrm{prec} \geq 0.75$ and $\mathrm{lift} > 1.1$. Requiring lift
matters: precedence alone admits pairs separated only by the academic calendar, about half
the candidates. The resulting graph is acyclic without intervention.

\figslot{fig2.png}{Knowledge graph enhanced knowledge tracing}{fig:kt}

\subsection{Knowledge Graph-Enhanced Knowledge Tracing}

Student knowledge concepts are estimated through past encounters by adopting a knowledge
tracing model that is based on a Transformer \cite{vaswani2017attention}. Temporal
attention depicts the dependency of learning activities and changes in learners' knowledge
over time. A Graph Convolutional Network (GCN) \cite{kipf2017gcn} simultaneously encodes
the precondition relationships in the knowledge graph, producing a concept embedding that
is recomputed at every forward pass so that gradients propagate into the graph structure.
The combined temporal learner representation and graph-based concept representation make
up the current learner state. The system may take into account the learner's past progress
as well as the structural links between ideas thanks to this integrated representation.

Forgetting is modelled explicitly, as an additive attention bias that decays with elapsed
time,
%
\begin{equation}
b^{(h)}_{ij} = -\,\mathrm{softplus}(\gamma_h)\,\log(1 + \Delta_{ij}),
\label{eq:decay}
\end{equation}
%
where $\Delta_{ij}$ is the number of days between interactions $i$ and $j$ and $\gamma_h$
is learned separately for each attention head, so that one head may attend to recent
activity while another retains longer memory. Setting $\gamma_h = 0$ recovers ordinary
attention and forms the temporal ablation reported in Section~\ref{sec:results}.

Two signals drive the model. Clickstream interactions update the hidden state but
contribute no training loss; assessment interactions are the supervised positions, and
binary cross-entropy is computed over those alone. This dual-signal construction is what
makes a sequence model of length 200 trainable on a dataset in which learners average
$8.7$ assessments each. Because predicted mastery is displayed to learners inside
explanations, the output probabilities are calibrated by temperature scaling
\cite{guo2017calibration}, with the temperature fitted on one half of each fold's
held-out predictions and evaluated on the other.

\figslot{fig3.png}{Prerequisite-constrained learning path optimization}{fig:planner}

\subsection{Prerequisite-Constrained Path Optimization}

The recommendation problem is formulated as a Markov Decision Process, where each course,
module or activity is a course of action, and the learner state is the knowledge and
learning context of the learner. The state concatenates the estimated mastery vector over
all concepts, a coverage mask of concepts already recommended, the static learner features
and the remaining budget. The reward combines learning progress, prerequisite
satisfaction, a redundancy penalty and a small step cost.

Rather than expecting a learned policy to discover the pedagogical constraint, it is
encoded directly in the set of admissible actions: a concept is admissible only if the
learner has not already mastered it, every prerequisite is at or above a readiness
threshold, and it belongs to an enrolled module. An unprepared recommendation is therefore
impossible by construction rather than merely unlikely, as Section~\ref{sec:results}
verifies.

The policy evaluated here is greedy with respect to one-step predicted improvement: each
admissible concept is simulated as studied, the resulting mastery vector is compared
against a baseline in which the same period elapses without study, and the concept giving
the largest total improvement in log-odds is recommended. Two details are load-bearing.
The baseline must be an idle period rather than the present state, because studying
advances the clock and the decay of \eqref{eq:decay} then attenuates every earlier
interaction, so measuring against the present moment makes every action appear harmful.
The comparison must be in log-odds, because engaged learners saturate near $0.98$ where
probability differences round to zero and the ranking degenerates. Mastery is bimodal, so
the mastery and readiness thresholds are the 70th and 30th percentiles of each learner's
own distribution rather than fixed values. A learned policy over the same state and reward
is left to future work; Section~\ref{sec:results} reports the headroom for one.

\figslot{fig4.png}{Explainable AI module for recommendation justification}{fig:xai}

\subsection{Explainable Recommendation and Evaluation}

Once an activity has been chosen, the priority-backtracking interpreter follows its choice
across the knowledge graph to find pertinent mastery levels, prerequisite links, and
knowledge gaps. Each prerequisite ancestor is scored by its shortfall below the readiness
threshold, discounted by graph distance, so that an unmet prerequisite immediately
preceding the recommendation outranks an equal shortfall further back; ancestors are
reported in descending order of that score. Counterfactual statements are computed by
raising a gap concept's mastery to the mastery threshold and re-evaluating the planner's
score, so the claimed effect of closing a gap is measured rather than asserted. These
elements are then turned into pedagogically meaningful explanations for students and
academic advisors \cite{zhang2023finegrained}.

AUC-ROC, RMSE and expected calibration error are used to assess knowledge-tracing
performance. Learning-path quality is assessed against observed behaviour rather than
simulation: each held-out learner's history is truncated at several points, each planner
is asked for a recommendation, and the ground truth is the set of concepts that learner
actually engaged with over the following interactions, with only learners who passed or
achieved distinction contributing ground truth. Explanation quality is assessed by three
objective measures --- fidelity, sufficiency and stability --- defined below, rather than by participant rating. Baselines comprise a
majority-class predictor, SAKT, DKT and a GNN-based recommender. Ablation experiments are
performed to remove certain elements from the proposed framework to determine the effect
on overall performance, under both the conventional protocol and the cold-concept protocol
introduced in Section~\ref{sec:results}.

% =====================================================================
\section{Implementation}
\label{sec:impl}

\subsection{Development Tools and Computing Environment}

The framework is implemented in Python 3.11. The Transformer-based knowledge tracing model
and the graph convolution component use PyTorch. Data preparation is performed in SQL
against an embedded DuckDB database, so every transformation from raw records to model
input is auditable without reading application code. NetworkX builds and traverses the
knowledge graph, and the \texttt{mlxtend} implementation of FP-Growth mines co-requisites.
Training used a single NVIDIA T4 GPU; all data processing, graph construction, planning
and inference run on CPU, on an Intel Core i7 workstation with 16~GB of RAM. Data
preparation takes about 2.5 minutes and training one configuration about 25 minutes.

\subsection{Dataset Configuration and Preprocessing}

The main dataset for implementation and assessment is the Open University Learning
Analytics Dataset (OULAD). The data set includes clickstream interaction data, student
enrollment, assessment scores and student demographics, comprising 32{,}593 enrolments,
10{,}655{,}280 clickstream events and 173{,}912 assessment records across seven modules.
A pipeline built on SQL, Pandas and NumPy preprocesses the raw records into sequential
learning data. Interactions are presented in chronological order and grouped into
learner-days. The categorical variables such as gender and education level are converted
to one-hot variables and the missing assessment scores are handled using median
imputation, applied only where a submission exists.

Two properties of the data required specific treatment and are reported here because both
materially affect what the reported metrics mean.

\textbf{Label construction.} OULAD records only \emph{submitted} assessments. Binarising
at the pass mark of 40 yields a $95.6\%$ positive rate, the fifth percentile of submitted
scores being the pass mark itself: the failures are absent from the records rather than
from reality. Of 323{,}925 expected submissions, 150{,}013 never occurred. Treating
non-submission as failure, restricted to assessments falling while a learner remained
registered, gives a $67.9/32.1$ balance and makes the target successful completion.

\textbf{Concept supervision is sparse and uneven.} The implementation contains 237 concept
nodes, 724 prerequisite relationships and 302 co-requisite relationships, with no cycles
requiring removal. Assessments exist for only 89 of the 237 concepts; mastery estimates
for the remaining 148 are reached through prerequisite propagation. This is the condition
that Section~\ref{sec:results} examines directly. Temporal information such as the volume
of study activity, elapsed time since the previous interaction and progression in mastery
of concepts is added to represent changes in learner behaviour and knowledge over time.

\subsection{Model Implementation Specifications}

A Transformer-based knowledge tracing encoder with two hidden layers, a hidden dimension
of 128, four attention heads, and a maximum sequence length of 200 interactions is used in
the learner modeling component. The three graph convolution layers in the graph
convolution (GCN) component are followed by layer normalization and ReLU activation.
64-dimensional concept embeddings are produced by the resulting graph encoder. These
embeddings are combined with the temporal learner representation to form the state on
which the planner operates. Dropout is $0.2$. Optimisation uses Adam at learning rate
$10^{-3}$ with cosine annealing, weight decay $10^{-4}$, batch size 64, and at most 12
epochs with early stopping at patience 3. The complete model has 466{,}569 parameters.
Random seeds are fixed and recorded alongside every result.

The planner evaluates all admissible concepts in one batched forward pass, which keeps
recommendation latency at roughly 160~ms on CPU. The priority-backtracking interpreter keeps track of the
relationship between a suggested idea and the concepts that it requires by means of
recursive graph traversal to a maximum depth of four. It draws out pedagogically
meaningful explanations for each recommendation, based on the learner proficiency levels
and relevant gaps in knowledge.

\subsection{Evaluation and Validation Implementation}

To check the implemented system against the majority-class, SAKT, DKT and GNN baselines, a
consistent evaluation technique is used. To evaluate predictive performance, AUC-ROC, RMSE
and expected calibration error are employed, while learning-path quality is measured by
NDCG@5, Hit@5, Precision@5, prerequisite violation rate and concept coverage against the
concepts learners actually studied next. Explanation quality is measured by fidelity,
sufficiency and stability. Ablation tests are performed to remove certain elements from
the proposed framework to determine the effect on overall performance. Five-fold cross-validation stratified on final outcome and grouped by
student is used, so that no learner contributes to both training and test partitions, with
paired $t$-tests over folds and Holm--Bonferroni correction \cite{holm1979} across the
comparison family. Cohen's $d$ is reported alongside every $p$-value, since with large
numbers of paired decisions significance is attainable at negligible effect size.

% =====================================================================
\section{Results and Discussion}
\label{sec:results}

\subsection{Predictive Accuracy Performance}

The proposed framework was evaluated against four baselines: a majority-class predictor,
SAKT, DKT and a GNN-based recommender. As shown in Table~\ref{tab:accuracy}, the suggested
framework achieves the highest AUC-ROC at $0.9538 \pm 0.0019$ and the lowest RMSE at
$0.2621 \pm 0.0026$, exceeding DKT by $0.0103$ ($p = 0.0009$) and the best baseline, the
GNN-based model, by $0.0036$ ($p = 0.022$) after Holm--Bonferroni correction. The results
indicate that temporal knowledge tracing and graph-based concept links can be combined to
get a more accurate representation of the changing knowledge state of the learner.

Two controls support that reading. The majority-class baseline returns exactly $0.5000$,
confirming the discrimination is not an artefact of the $67.9/32.1$ class distribution,
and SAKT --- attention without graph conditioning or temporal decay --- reaches $0.9233$,
so the improvement does not follow simply from using attention. Expected calibration error
falls from $0.0084$ to $0.0054$ after temperature scaling; since predicted mastery is
shown to learners inside explanations, a poorly calibrated probability is an incorrect
statement rather than a cosmetic defect.

\subsection{Learning Path Quality}

The learning-path results show how successful sequential optimization is.
Table~\ref{tab:pathquality} reports 906 recommendation decisions drawn from 400 held-out
learners who passed or achieved distinction, with the concepts each learner actually
studied next as ground truth. The proposed framework achieves the highest NDCG@5 at
$0.1532$ and the highest Hit@5 at $0.3687$, ahead of weakest-first ($0.1344$), random
selection among admissible actions ($0.1300$), popularity ($0.1069$) and syllabus order
($0.0612$). All differences are significant after correction.

The effect sizes qualify that conclusion, and we report them rather than the $p$-values
alone. Against syllabus order the effect is small ($d = 0.36$); against random selection
among admissible actions it is negligible ($d = 0.09$). The honest interpretation is that
the prerequisite constraint performs most of the work and the learned ranking adds a small
but consistent improvement on top of it. This also bounds the headroom available to a
learned policy: an agent optimising the reward defined in Section~\ref{sec:method} would
have to improve on a greedy baseline that itself exceeds random selection by only
$d = 0.09$.

Prerequisite violations are zero across all 906 decisions and every planner, so the
guarantee of Section~\ref{sec:method} is verified rather than assumed. Coverage of $0.726$
shows the planner does not collapse onto a few popular concepts, the failure mode
popularity-based recommendation exhibits here at $0.510$.

\subsection{Explanation Quality}

The explainability module was evaluated by three objective measures rather than by
participant rating, over 25 recommendation decisions each repeated under three
perturbations. The sample is small and reported as such, but each measure is computed
against the model's own counterfactual response rather than against an opinion. Fidelity of $0.5318$ indicates that cited prerequisite deficits are positively but
imperfectly predictive of measured counterfactual improvement. Stability of $0.5310$
indicates that roughly half the cited gaps persist under a $\pm 0.02$ perturbation of the
mastery estimate.

Sufficiency of $0.0020$ is the most informative of the three and is reported here as a
negative finding. It shows that the concept named in an explanation accounted for only
$0.2\%$ of the predicted benefit, with the remainder arising from transfer to related
concepts that the explanation did not mention. The explanations were accurate but
substantially incomplete, and the renderer was subsequently amended to name the principal
beneficiaries. No rating scale would have surfaced this, because participants cannot
assess a quantity that the explanation omits.

\subsection{Ablation Study}

The lower rows of Table~\ref{tab:accuracy} report the conventional ablation. Removing knowledge-graph encoding changes predictive accuracy by
$-0.0001$ ($p = 0.77$), and removing the temporal attention bias changes it by $-0.0006$
($p = 0.45$). Read at face value, neither component contributes anything, which
contradicts the expectation that knowledge-graph encoding supplies important structural
information.

\subsection{Why Conventional Ablation Is Insufficient}

That reading is an artefact of the measurement rather than a property of the model.
Predictive accuracy is computed only at assessment positions, and assessments exist for 89
of the 237 concepts, each carrying on the order of 2{,}700 labelled examples. A free
embedding table learns such concepts adequately without any structural information, so the
graph is genuinely redundant there. Its function is to support the remaining 148 concepts,
which never enter the accuracy computation at all. The conventional ablation is
structurally incapable of observing the component it removes.

To measure the contribution directly, all assessments were withheld for 27 of the 89
assessed concepts and their observed responses were removed from the model input, so that
those concepts behave as genuinely unassessed; evaluation was then restricted to them
alone. Table~\ref{tab:coldconcept} reports the outcome: $0.9004 \pm 0.0059$ with graph
propagation against $0.8679 \pm 0.0042$ without, a difference of $+0.0325$
($p = 0.0003$, Cohen's $d = 5.06$), with the graph-conditioned model superior on every
fold.

Knowledge-graph propagation therefore contributes nothing where direct supervision is
plentiful and substantially where it is absent, and a conventional ablation cannot
distinguish the two cases. This has a practical consequence for the recommender, which
must rank across all 237 concepts, most of which carry no assessment of their own; and a
methodological consequence for the wider literature, in which component ablations are
routinely reported on metrics that are unable to observe the component under test. We
suggest that such ablations be accompanied by evidence that the reporting metric is
capable of detecting the component being removed.

\subsection{Overall Discussion}

The experimental results demonstrate that combining explainability, knowledge-graph
modeling, constrained path optimization and knowledge tracing improves both predictive and
recommendation-oriented metrics relative to the assessed baselines. The explainability
process enhances the interpretability of suggestions without being directly in charge of
forecast accuracy. It does increase computational cost: the average time required to
produce a fully explained recommendation is approximately 160~ms on CPU, against 7~ms for
a cached mastery lookup alone. That trade-off indicates that the proposed framework
provides substantially more transparency at an acceptable computational cost, and remains
well within the latency budget of an interactive advising interface.

% =====================================================================
\section{Conclusion and Future Work}
\label{sec:conclusion}

\subsection{Conclusion}

This study employs Knowledge Graph modeling, pedagogical explanation generation,
prerequisite-constrained path optimization and Deep Knowledge Tracing to develop an
Explainable AI based personalized route recommendation system for college students. The
framework builds optimized learning paths by modelling learners' changing knowledge states
together with precondition links and temporal behavior, and the priority-backtracking
process finds the knowledge gaps and dependencies that justify each suggestion.

On the Open University Learning Analytics Dataset the framework exceeded the
majority-class, SAKT, DKT and GNN-based baselines, with an AUC-ROC of
$0.9538 \pm 0.0019$, an RMSE of $0.2621 \pm 0.0026$ and a calibration error of $0.0054$.
The planner produced zero prerequisite violations across 906 decisions evaluated against
real learner behaviour, with the qualification that the prerequisite constraint rather than
the learned ranking accounts for most of the margin over random selection.

The principal finding of the ablation study concerns evaluation. A conventional ablation
indicates that knowledge-graph propagation contributes nothing, yet under a protocol that
evaluates concepts lacking direct supervision the same component improves accuracy by
$+0.0325$ with $d = 5.06$. The discrepancy is a property of the metric, not of the model:
aggregate accuracy measured where supervision is abundant cannot observe a structural
prior whose purpose is to compensate for its absence. The results indicate that learner-state modeling, structural
knowledge representation, constrained sequence optimization and explanation generation
together provide a transparent and personalized learning-path recommendation strategy for
higher education --- and that verifying such a claim requires evaluation protocols chosen
to match what each component is for.

\subsection{Limitations}

The prediction target is successful completion rather than knowledge. Because OULAD records
only submitted assessments, treating non-submission as failure is necessary for a usable
class balance, but a substantial share of the signal is then engagement rather than
understanding --- appropriate for a system intended to prompt intervention, but a weaker
claim than measuring mastery. The concept graph is likewise derived rather than given:
OULAD supplies no concept annotations, prerequisite relations or material titles, so
concepts are defined structurally and labelled by activity composition. Placement for
$82\%$ of materials relies on observed access, validated at $\rho = 0.975$ on the subset
carrying both signals, but deriving concept order from behaviour and then mining
prerequisites from behaviour is partially self-referential.

Supervision covers 89 of 237 concepts, so mastery for the remainder is not directly
calibrated, and predicted mastery is bimodal, which is why learner-relative thresholds are
required. The proposed model additionally conditions on learner features that the DKT and
SAKT baselines do not, as those follow their original formulations; part of the reported
improvement may therefore reflect feature access rather than architecture. Finally, all
results derive from a single institution's distance-learning data across seven modules,
and generalisation to campus-based teaching is untested.

\subsection{Ethical Considerations}

All data used in this study is drawn from OULAD, which is publicly released under CC-BY
4.0 and contains no personally identifying information. No data was collected from human
participants. The system produces advisory recommendations and makes no determinative
decision concerning any learner. It does, however, condition on attributes including
gender, disability status, index of multiple deprivation and age band, and a subgroup
analysis across these attributes has not been conducted. Until one is, the system should
not be deployed in any setting where its recommendations influence assessment or
progression decisions.

\subsection{Future Work}

Four directions follow from the limitations above. The most immediate is a \textbf{fairness
analysis}: subgroup performance across gender, disability, deprivation band and age band
must be established before any deployment. Second, a \textbf{learned policy} over the state
and reward of Section~\ref{sec:method} should be trained and evaluated, in simulation and
by off-policy estimation on logged trajectories; Section~\ref{sec:results} quantifies the
headroom, since the greedy planner exceeds random selection among admissible actions by
only $d = 0.09$, and algorithms such as Soft Actor-Critic or Proximal Policy Optimization
would have to widen that margin to justify their complexity. Third, a \textbf{user
evaluation} would address what the objective measures cannot: fidelity, sufficiency and
stability establish whether explanations are faithful to the model, not whether learners
find them useful. Fourth, \textbf{expert validation of the mined edges} would give an
independent estimate of graph precision, since the prerequisite relations are inferred
from behaviour rather than supplied by a curriculum designer.

Beyond these, future studies can integrate multimodal information such as forum
participation and video-lecture interaction, examine transfer across subjects and
institutions, and create lightweight versions through pruning and distillation for
deployment on resource-limited platforms.

\subsection*{Reproducibility}

OULAD is available under CC-BY 4.0. All source code, SQL transformations and generated
results are at \texttt{github.com/abhay27singh/\allowbreak explainable-learning-path-recommender};
every table here is regenerated from raw data by one command with fixed seeds.

% =====================================================================
%                              TABLES
% =====================================================================

\begin{table}[t]
\centering
\caption{Comparative Predictive Performance and Component Ablation}
\label{tab:accuracy}
\setlength{\tabcolsep}{3pt}\footnotesize
\begin{tabular}{lccc}
\hline
\textbf{Model} & \textbf{AUC-ROC} & \textbf{RMSE} & \textbf{$p$} \\
\hline
Majority (no-skill) & $0.5000 \pm 0.0000$ & $0.4669$ & $<0.001$ \\
SAKT                & $0.9233 \pm 0.0018$ & $0.2992$ & $<0.001$ \\
DKT                 & $0.9435 \pm 0.0017$ & $0.2751$ & $0.0009$ \\
GNN-based           & $0.9502 \pm 0.0014$ & $0.2665$ & $0.0220$ \\
\textbf{Proposed}   & $\mathbf{0.9538 \pm 0.0019}$ & $\mathbf{0.2621}$ & --- \\
\hline
\quad w/o knowledge graph    & $0.9539 \pm 0.0021$ & $0.2618$ & $0.7704$ \\
\quad w/o temporal attention & $0.9532 \pm 0.0014$ & $0.2624$ & $0.4507$ \\
\hline
\end{tabular}

\vspace{2pt}
{\footnotesize Folds grouped by student; $p$ Holm-corrected against the proposed model.
Proposed ECE $=0.0054$. Neither ablated component shows a detectable change --- see
Section~\ref{sec:results}-E.}
\end{table}

\begin{table}[t]
\centering
\caption{Learning Path Quality Comparison}
\label{tab:pathquality}
\begin{tabular}{lcccc}
\hline
\textbf{Planner} & \textbf{NDCG@5} & \textbf{Hit@5} & \textbf{Viol.} & \textbf{Cov.} \\
\hline
Syllabus order      & $0.0612$ & $0.1832$ & $0.0000$ & $0.452$ \\
Popularity          & $0.1069$ & $0.2936$ & $0.0000$ & $0.510$ \\
Random (admissible) & $0.1300$ & $0.3455$ & $0.0000$ & $0.722$ \\
Weakest-first       & $0.1344$ & $0.3642$ & $0.0000$ & $0.515$ \\
\textbf{Proposed}   & $\mathbf{0.1532}$ & $\mathbf{0.3687}$ & $\mathbf{0.0000}$ & $\mathbf{0.726}$ \\
\hline
\end{tabular}

\vspace{2pt}
{\footnotesize 906 decisions, 400 held-out successful learners; ground truth is the
concepts each actually studied next. NDCG@5 vs.\ proposed: syllabus $+0.0921$ ($d=0.36$);
popularity $+0.0464$ ($d=0.17$); random $+0.0232$ ($d=0.09$); weakest-first $+0.0189$
($d=0.08$); all $p<0.02$ Holm-corrected.}
\end{table}

\begin{table}[t]
\centering
\caption{Cold-Concept Ablation}
\label{tab:coldconcept}
\begin{tabular}{lcc}
\hline
\textbf{Variant} & \textbf{AUC-ROC} & \textbf{RMSE} \\
\hline
With knowledge graph & $\mathbf{0.9004 \pm 0.0059}$ & $\mathbf{0.3359}$ \\
Without              & $0.8679 \pm 0.0042$ & $0.3611$ \\
\hline
\textbf{Difference}  & \multicolumn{2}{c}{$\mathbf{+0.0325}$, $p = 0.0003$, $d = 5.06$} \\
\hline
\end{tabular}

\vspace{2pt}
{\footnotesize All assessments withheld for 27 of the 89 assessed concepts, their responses
stripped from the input, evaluation restricted to those concepts. Graph wins every fold.}
\end{table}

% =====================================================================
\begin{thebibliography}{00}

\bibitem{li2026survey}
A.~Li, Y.~Li, and X.~Gao, ``Personalized learning path recommendation based on knowledge
graphs: A survey,'' \emph{Electronics}, vol.~15, no.~1, Art. no. 238, 2026, doi:
10.3390/electronics15010238.

\bibitem{aher2013combination}
S.~B. Aher and L.~M.~R.~J. Lobo, ``Combination of machine learning algorithms for
recommendation of courses in e-learning system based on historical data,''
\emph{Knowledge-Based Systems}, vol.~51, pp.~1--14, 2013, doi:
10.1016/j.knosys.2013.04.015.

\bibitem{kuzilek2017oulad}
J.~Kuzilek, M.~Hlosta, and Z.~Zdrahal, ``Open University Learning Analytics dataset,''
\emph{Scientific Data}, vol.~4, Art. no. 170171, 2017, doi: 10.1038/sdata.2017.171.

\bibitem{corbett1994knowledge}
A.~T. Corbett and J.~R. Anderson, ``Knowledge tracing: Modeling the acquisition of
procedural knowledge,'' \emph{User Modeling and User-Adapted Interaction}, vol.~4, no.~4,
pp.~253--278, 1994.

\bibitem{piech2015dkt}
C.~Piech, J.~Bassen, J.~Huang, S.~Ganguli, M.~Sahami, L.~Guibas, and J.~Sohl-Dickstein,
``Deep knowledge tracing,'' in \emph{Advances in Neural Information Processing Systems},
vol.~28, 2015, pp.~505--513.

\bibitem{zhang2017dkvmn}
J.~Zhang, X.~Shi, I.~King, and D.-Y. Yeung, ``Dynamic key-value memory networks for
knowledge tracing,'' in \emph{Proc. 26th Int. World Wide Web Conf. (WWW)}, 2017,
pp.~765--774, doi: 10.1145/3038912.3052580.

\bibitem{ghosh2020akt}
A.~Ghosh, N.~T. Heffernan, and A.~S. Lan, ``Context-aware attentive knowledge tracing,''
in \emph{Proc. 26th ACM SIGKDD Int. Conf. Knowledge Discovery and Data Mining (KDD)},
2020, pp.~2330--2339, doi: 10.1145/3394486.3403282.

\bibitem{shen2021lpkt}
S.~Shen, Q.~Liu, E.~Chen, Z.~Huang, W.~Huang, Y.~Yin, Y.~Su, and S.~Wang, ``Learning
process-consistent knowledge tracing,'' in \emph{Proc. 27th ACM SIGKDD Conf. Knowledge
Discovery and Data Mining (KDD)}, 2021, pp.~1452--1460, doi: 10.1145/3447548.3467237.

\bibitem{vaswani2017attention}
A.~Vaswani, N.~Shazeer, N.~Parmar, J.~Uszkoreit, L.~Jones, A.~N. Gomez, L.~Kaiser, and
I.~Polosukhin, ``Attention is all you need,'' in \emph{Advances in Neural Information
Processing Systems}, vol.~30, 2017, pp.~5998--6008.

\bibitem{kipf2017gcn}
T.~N. Kipf and M.~Welling, ``Semi-supervised classification with graph convolutional
networks,'' in \emph{Proc. Int. Conf. Learning Representations (ICLR)}, 2017.

\bibitem{mnih2015dqn}
V.~Mnih \emph{et al.}, ``Human-level control through deep reinforcement learning,''
\emph{Nature}, vol.~518, pp.~529--533, 2015, doi: 10.1038/nature14236.

\bibitem{khosravi2022xaied}
H.~Khosravi, S.~B. Shum, G.~Chen, C.~Conati, Y.-S. Tsai, J.~Kay, S.~Knight,
R.~Martinez-Maldonado, S.~Sadiq, and D.~Ga\v{s}evi\'{c}, ``Explainable artificial
intelligence in education,'' \emph{Computers and Education: Artificial Intelligence},
vol.~3, Art. no. 100074, 2022, doi: 10.1016/j.caeai.2022.100074.

\bibitem{ribeiro2016lime}
M.~T. Ribeiro, S.~Singh, and C.~Guestrin, ``Why should I trust you? Explaining the
predictions of any classifier,'' in \emph{Proc. 22nd ACM SIGKDD Int. Conf. Knowledge
Discovery and Data Mining}, 2016, pp.~1135--1144, doi: 10.1145/2939672.2939778.

\bibitem{lundberg2017shap}
S.~M. Lundberg and S.-I. Lee, ``A unified approach to interpreting model predictions,''
in \emph{Advances in Neural Information Processing Systems}, vol.~30, 2017,
pp.~4765--4774.

\bibitem{zhang2023finegrained}
S.~Zhang, N.~Hui, P.~Zhai, J.~Xu, L.~Cao, and Q.~Wang, ``A fine-grained and
multi-context-aware learning path recommendation model over knowledge graphs for online
learning communities,'' \emph{Information Processing \& Management}, vol.~60, no.~5, Art.
no. 103464, 2023, doi: 10.1016/j.ipm.2023.103464.

\bibitem{yang2025explainable}
Y.~Yang, X.~Peng, M.~Chen, and S.~Liu, ``An explainable graph-based course recommendation
model based on multiple interest factors,'' \emph{Expert Systems with Applications},
vol.~264, Art. no. 125889, 2025, doi: 10.1016/j.eswa.2024.125889.

\bibitem{xie2025profile}
X.~Xie and X.~Feng, ``Personalized learning path recommendation based on learner profile
and knowledge graph,'' \emph{Computing and Informatics}, vol.~44, no.~4, pp.~983--1008,
2025, doi: 10.31577/cai-202-4983.

\bibitem{zhao2026rlkt}
J.~Zhao, J.~Kuang, S.~Deng, R.~Kuang, and Z.~Zhang, ``Personalized learning path
recommendation for middle school students based on joint optimization of deep
reinforcement learning and knowledge tracing,'' \emph{Scientific Reports}, 2026, doi:
10.1038/s41598-026-61166-6.

\bibitem{boujmiraz2026review}
S.~Boujmiraz, H.~Darhmaoui, and A.~Drissi El Maliani, ``Predicting student performance: A
comprehensive review of machine learning, deep learning, and explainable AI approaches,''
\emph{Computers and Education: Artificial Intelligence}, vol.~10, Art. no. 100548, 2026,
doi: 10.1016/j.caeai.2026.100548.

\bibitem{qiang2026xaiedu}
M.~Qiang, Z.~Liu, and R.~Zhang, ``Explainable AI in education: Integrating educational
domain knowledge into the deep learning model for improved student performance
prediction,'' \emph{Scientific Reports}, vol.~16, Art. no. 9515, 2026, doi:
10.1038/s41598-026-40538-y.

\bibitem{pandey2019sakt}
S.~Pandey and G.~Karypis, ``A self-attentive model for knowledge tracing,'' in
\emph{Proc. 12th Int. Conf. Educational Data Mining (EDM)}, 2019, pp.~384--389.

\bibitem{guo2017calibration}
C.~Guo, G.~Pleiss, Y.~Sun, and K.~Q. Weinberger, ``On calibration of modern neural
networks,'' in \emph{Proc. 34th Int. Conf. Machine Learning (ICML)}, 2017,
pp.~1321--1330.

\bibitem{holm1979}
S.~Holm, ``A simple sequentially rejective multiple test procedure,''
\emph{Scandinavian Journal of Statistics}, vol.~6, no.~2, pp.~65--70, 1979.

\end{thebibliography}

\end{document}
